\renewcommand{\abstractname}{Resumo}

Este trabalho propõe um chatbot baseado em regras para auxiliar discentes, docentes
e o público acadêmico em geral na busca de atas institucionais da Universidade
Federal Rural do Rio de Janeiro (UFRRJ). O crescimento contínuo do volume de
atas torna a busca manual cada vez mais custosa e para resolver esse problema,
propõe-se um roteiro guiado de diálogo que conduz o usuário por um conjunto fechado
de critérios de busca, sendo eles: departamento, disciplina, docentes, discentes, período,
processo/edital e progressão funcional, além de um campo de assuntos gerais,
traduzindo cada resposta em um parâmetro estruturado de uma consulta ao
Elasticsearch. A solução foi implementada na plataforma Botpress v12, integrada
a um indexador desenvolvido em Node.js responsável por extrair o conteúdo de
arquivos PDF e montar dinamicamente a consulta em formato Query DSL, além de um bot para o Discord como canal alternativo de interação. Os documentos
retornados são ranqueados pelo algoritmo BM25 e são apresentadas ao usuário as 3 atas com maior relevância. Os resultados indicam que o roteiro guiado de diálogo constitui uma solução
viável para a geração de consultas estruturadas a partir de interações em linguagem
natural, deixando claro o que pode ou não ser buscado e dispensando
do usuário qualquer conhecimento sobre a organização do índice ou a sintaxe de
consulta do motor de busca.

\textbf{Palavras-chave:} chatbot baseado em regras; recuperação da informação;
Elasticsearch; atas institucionais; Botpress.

